\documentclass[10pt,a4paper]{amsart}
\usepackage{amsmath,amssymb,amsthm,graphicx}
% Retained preamble for the Stacks Project chapter snapshot in this directory.
% The chapter's own preamble.tex is a copy of the shared Stacks Project preamble; the block
% below is sliced MECHANICALLY, line for line, from it. Nothing was added, removed or
% redefined; the wrapper class, the xr-hyper/hyperref wiring and the externaldocument list
% are omitted because the assembler supplies the document class.
%
% lines taken from preamble.tex: 3, 5, 9, 11, 12, 156, 157, 158, ... (51 lines total)
%
% package coverage restored from the chapter's own preamble (2026-10-13): xy, xy-2cell

\usepackage[all]{xy}
\xyoption{2cell}
\UseAllTwocells

\newenvironment{reference}{\comment}{\endcomment}
\newenvironment{slogan}{\comment}{\endcomment}
\newenvironment{history}{\comment}{\endcomment}
\theoremstyle{plain}
\newtheorem{theorem}[subsection]{Theorem}
\newtheorem{proposition}[subsection]{Proposition}
\newtheorem{lemma}[subsection]{Lemma}
\theoremstyle{definition}
\newtheorem{definition}[subsection]{Definition}
\newtheorem{example}[subsection]{Example}
\newtheorem{exercise}[subsection]{Exercise}
\newtheorem{situation}[subsection]{Situation}
\theoremstyle{remark}
\newtheorem{remark}[subsection]{Remark}
\newtheorem{remarks}[subsection]{Remarks}
\numberwithin{equation}{subsection}
\def\lim{\mathop{\mathrm{lim}}\nolimits}
\def\colim{\mathop{\mathrm{colim}}\nolimits}
\def\Spec{\mathop{\mathrm{Spec}}}
\def\Hom{\mathop{\mathrm{Hom}}\nolimits}
\def\Ext{\mathop{\mathrm{Ext}}\nolimits}
\def\SheafHom{\mathop{\mathcal{H}\!\mathit{om}}\nolimits}
\def\SheafExt{\mathop{\mathcal{E}\!\mathit{xt}}\nolimits}
\def\Sch{\mathit{Sch}}
\def\Mor{\mathop{\mathrm{Mor}}\nolimits}
\def\Ob{\mathop{\mathrm{Ob}}\nolimits}
\def\Sh{\mathop{\mathit{Sh}}\nolimits}
\def\NL{\mathop{N\!L}\nolimits}
\def\CH{\mathop{\mathrm{CH}}\nolimits}
\def\proetale{{pro\text{-}\acute{e}tale}}
\def\etale{{\acute{e}tale}}
\def\QCoh{\mathit{QCoh}}
\def\Ker{\mathop{\mathrm{Ker}}}
\def\Im{\mathop{\mathrm{Im}}}
\def\Coker{\mathop{\mathrm{Coker}}}
\def\Coim{\mathop{\mathrm{Coim}}}
\DeclareMathSymbol{\boxtimes}{\mathbin}{AMSa}{"02}
\def\QCohstack{\mathcal{QC}\!\mathit{oh}}
\def\Cohstack{\mathcal{C}\!\mathit{oh}}
\def\Spacesstack{\mathcal{S}\!\mathit{paces}}
\def\Quotfunctor{\mathrm{Quot}}
\def\Hilbfunctor{\mathrm{Hilb}}
\def\Curvesstack{\mathcal{C}\!\mathit{urves}}
\def\Polarizedstack{\mathcal{P}\!\mathit{olarized}}
\def\Complexesstack{\mathcal{C}\!\mathit{omplexes}}
\def\Pic{\mathop{\mathrm{Pic}}\nolimits}
\def\Picardstack{\mathcal{P}\!\mathit{ic}}
\def\Picardfunctor{\mathrm{Pic}}
\def\Deformationcategory{\mathcal{D}\!\mathit{ef}}
\renewenvironment{reference}{}{}
\renewenvironment{slogan}{}{}
\renewenvironment{history}{}{}
\usepackage{cleveref}
\crefname{originallemma}{original lemma}{original lemmas}
\Crefname{originallemma}{Original lemma}{Original lemmas}
\makeatletter
\newlabel{lemma-derived-completion-koszul}{{0A0E}{source}{}{sourcepointer-1}{}}
\newlabel{lemma-derived-completion-koszul@cref}{{[originallemma][0][]0A0E}{[1][0][]source}{}{}{}}
\makeatother
\crefname{originallemma}{original lemma}{original lemmas}
\Crefname{originallemma}{Original lemma}{Original lemmas}
\makeatletter
\newlabel{lemma-pushforward-derived-complete}{{099J}{source}{}{sourcepointer-2}{}}
\newlabel{lemma-pushforward-derived-complete@cref}{{[originallemma][0][]099J}{[1][0][]source}{}{}{}}
\makeatother
\crefname{originallemma}{original lemma}{original lemmas}
\Crefname{originallemma}{Original lemma}{Original lemmas}
\makeatletter
\newlabel{proposition-derived-completion}{{099F}{source}{}{sourcepointer-3}{}}
\newlabel{proposition-derived-completion@cref}{{[originallemma][0][]099F}{[1][0][]source}{}{}{}}
\makeatother
\crefname{originallemma}{original lemma}{original lemmas}
\Crefname{originallemma}{Original lemma}{Original lemmas}
\makeatletter
\newlabel{remark-localization-and-completion}{{0A0F}{source}{}{sourcepointer-4}{}}
\newlabel{remark-localization-and-completion@cref}{{[originallemma][0][]0A0F}{[1][0][]source}{}{}{}}
\makeatother
\crefname{originallemma}{original lemma}{original lemmas}
\Crefname{originallemma}{Original lemma}{Original lemmas}
\makeatletter
\newlabel{sites-cohomology-lemma-Rf-commutes-with-Rlim}{{0A07}{source}{}{sourcepointer-5}{}}
\newlabel{sites-cohomology-lemma-Rf-commutes-with-Rlim@cref}{{[originallemma][0][]0A07}{[1][0][]source}{}{}{}}
\makeatother
\crefname{originallemma}{original lemma}{original lemmas}
\Crefname{originallemma}{Original lemma}{Original lemmas}
\makeatletter
\newlabel{sites-cohomology-lemma-projection-formula}{{0944}{source}{}{sourcepointer-6}{}}
\newlabel{sites-cohomology-lemma-projection-formula@cref}{{[originallemma][0][]0944}{[1][0][]source}{}{}{}}
\makeatother
\title{For morphisms of ringed topoi with finitely generated ideal, pushforward commutes with derived completion}
\author{The Stacks Project Authors}
\date{stacks-project:algebraization.tex:local-snapshot-20261012; GNU FDL 1.2 or later}
\begin{document}
\maketitle
\noindent\textit{Editorial note.} Statement span is lines 1723-1736 of sources/stacks-algebraization/pinned-original-source.tex; proof span is lines 1738-1774, the complete author proof (35 body lines).\par\medskip
\noindent\textit{Editorial note.} Navigation stubs: 6 cited-but-not-included labels are supplied as non-load-bearing pointers with their pinned tags where the snapshot tag index has them.\par\medskip
\noindent\textit{Editorial note.} extra\_packages is empty on purpose: the retained preamble already loads xy with options and a second declaration clashes in the assembler wrapper.\par\medskip
\noindent\textit{Editorial note.} Text is retained verbatim, including the author's own typos and placeholders.\par\medskip
\noindent\textit{Editorial note.} Pool algebraization wave36 (al02 bib-salvage batch). Spans computed from \textbackslash{}label (16.5/16.6).\par\medskip
\noindent\textit{Editorial note.} External \textbackslash{}cite keys salvaged with card-local bibliography.tex entries copied verbatim from the snapshot my.bib (precedent S3-morph-blow-up-ample-family); URL/ISSN/EPRINT fields dropped as the preamble has no url package.\par\medskip
\noindent\textit{Editorial note.} statement contains reference environment with real content -> patch-sd-envs.tex applied (sd02 lesson).\par\medskip
\noindent\textit{Allowed standard prerequisites (editorial).}\begin{itemize}
\item Explicit prerequisite boundary. See why-field; in-chapter lemmas cited by name become stub prereqs; external cites resolved via the card-local bibliography.tex fragment.
\end{itemize}\medskip
\ifdefined\hypertarget\hypertarget{sourcepointer-1}{}\fi
\noindent\textit{Source reference.} Original lemma 0A0E: sources/stacks-algebraization/pinned-original-source.tex, source lines 1308--1319. Cited by the retained statement or proof but not part of the collected spans; supplied only as a navigation pointer. The argument does not depend on its content.\par\medskip
\ifdefined\hypertarget\hypertarget{sourcepointer-2}{}\fi
\noindent\textit{Source reference.} Original lemma 099J: sources/stacks-algebraization/pinned-original-source.tex, source lines 1199--1206. Cited by the retained statement or proof but not part of the collected spans; supplied only as a navigation pointer. The argument does not depend on its content.\par\medskip
\ifdefined\hypertarget\hypertarget{sourcepointer-3}{}\fi
\noindent\textit{Source reference.} Original lemma 099F: sources/stacks-algebraization/pinned-original-source.tex, source lines 1515--1520. Cited by the retained statement or proof but not part of the collected spans; supplied only as a navigation pointer. The argument does not depend on its content.\par\medskip
\ifdefined\hypertarget\hypertarget{sourcepointer-4}{}\fi
\noindent\textit{Source reference.} Original lemma 0A0F: sources/stacks-algebraization/pinned-original-source.tex, source lines 1558--1576. Cited by the retained statement or proof but not part of the collected spans; supplied only as a navigation pointer. The argument does not depend on its content.\par\medskip
\ifdefined\hypertarget\hypertarget{sourcepointer-5}{}\fi
\noindent\textit{Source reference.} Original lemma 0A07: sources/stacks-algebraization/pinned-original-source.tex, source lines 1774--1774. Cited by the retained statement or proof but not part of the collected spans; supplied only as a navigation pointer. The argument does not depend on its content.\par\medskip
\ifdefined\hypertarget\hypertarget{sourcepointer-6}{}\fi
\noindent\textit{Source reference.} Original lemma 0944: sources/stacks-algebraization/pinned-original-source.tex, source lines 1774--1774. Cited by the retained statement or proof but not part of the collected spans; supplied only as a navigation pointer. The argument does not depend on its content.\par\medskip
\begin{lemma}
\label{lemma-pushforward-commutes-with-derived-completion}
\begin{reference}
Generalization of \cite[Lemma 6.5.9 (2)]{BS}. Compare with
\cite[Theorem 6.5]{HL-P} in the setting of quasi-coherent modules
and morphisms of (derived) algebraic stacks.
\end{reference}
Let $f : (\Sh(\mathcal{D}), \mathcal{O}') \to (\Sh(\mathcal{C}), \mathcal{O})$
be a morphism of ringed topoi. Let $\mathcal{I} \subset \mathcal{O}$
be a finite type sheaf of ideals. Let $\mathcal{I}' \subset \mathcal{O}'$
be the ideal generated by $f^\sharp(f^{-1}\mathcal{I})$.
Then $Rf_*$ commutes with derived completion, i.e.,
$Rf_*(K^\wedge) = (Rf_*K)^\wedge$.
\end{lemma}

\begin{proof}
By Proposition \ref{proposition-derived-completion} the derived completion
functors exist. By Lemma \ref{lemma-pushforward-derived-complete} the object
$Rf_*(K^\wedge)$ is derived complete, and hence we obtain a canonical map
$(Rf_*K)^\wedge \to Rf_*(K^\wedge)$ by the universal property of derived
completion. We may check this map is an isomorphism locally on $\mathcal{C}$.
Thus, since derived completion commutes with localization
(Remark \ref{remark-localization-and-completion}) we may assume
that $\mathcal{I}$ is generated by global sections $f_1, \ldots, f_r$.
Then $\mathcal{I}'$ is generated by $g_i = f^\sharp(f_i)$. By
Lemma \ref{lemma-derived-completion-koszul}
we have to prove that
$$
R\lim \left(
Rf_*K \otimes^\mathbf{L}_\mathcal{O} K(\mathcal{O}, f_1^n, \ldots, f_r^n)
\right)
=
Rf_*\left(
R\lim
K \otimes^\mathbf{L}_{\mathcal{O}'} K(\mathcal{O}', g_1^n, \ldots, g_r^n)
\right)
$$
Because $Rf_*$ commutes with $R\lim$
(Cohomology on Sites, Lemma
\ref{sites-cohomology-lemma-Rf-commutes-with-Rlim})
it suffices to prove that
$$
Rf_*K \otimes^\mathbf{L}_\mathcal{O} K(\mathcal{O}, f_1^n, \ldots, f_r^n) =
Rf_*\left(
K \otimes^\mathbf{L}_{\mathcal{O}'} K(\mathcal{O}', g_1^n, \ldots, g_r^n)
\right)
$$
This follows from the projection formula (Cohomology on Sites, Lemma
\ref{sites-cohomology-lemma-projection-formula}) and the fact that
$Lf^*K(\mathcal{O}, f_1^n, \ldots, f_r^n) =
K(\mathcal{O}', g_1^n, \ldots, g_r^n)$.
\end{proof}

\begin{thebibliography}{9}
\bibitem{BS}
Bhatt, Bhargav and Scholze, Peter.
\newblock \emph{The pro-\'etale topology for schemes}.
\newblock 2013. preprint.
\bibitem{HL-P}
Halpern-Leistner, Daniel and Preygel, Anatoly.
\newblock \emph{Mapping stacks and categorical notions of properness}.
\newblock 2014.
\end{thebibliography}
\end{document}
